\section{Global spectral blocks on a fixed Pell cat map}\label{sec:pell_blocks}
Unlike the separate magnetic-quotient semicircle theorem, this result uses a fixed smooth Fourier observable, deterministic parameters, and generally composite Pell moduli. No assertion that infinitely many Pell numbers are prime is needed. The two limiting laws concern different regimes and are not interchangeable.

\subsection{Arithmetic moduli and a fixed classical map}\label{pell_blocks:arithmetic-moduli-and-a-fixed-classical-map}

Put \[
 \alpha=1+\sqrt2,\qquad \lambda=\alpha^2=3+2\sqrt2,\qquad
 T=\begin{pmatrix}3&2\\4&3\end{pmatrix}.
\] The integral matrix T has determinant one and eigenvalues \(\lambda\) and its reciprocal, so it is one fixed hyperbolic toral automorphism. Also T is the identity modulo 2, avoiding the usual centered-Weyl parity ambiguity.

Define integers \(Q_{j}\),\(P_{j}\) by \[
 \alpha^j=Q_j+P_j\sqrt2,
 \quad Q_0=1,\ P_0=0,\quad
 Q_{j+1}=Q_j+2P_j,\ P_{j+1}=Q_j+P_j.
\] Let K tend to infinity through odd primes and set \[
 N=Q_K=\frac{\alpha^K-\alpha^{-K}}2,\qquad F=P_K.
\] Then \[
 N^2-2F^2=-1,\qquad N\asymp\lambda^{K/2}.
\] Every \(Q_{j}\) is odd, and \(\gcd(F,N)=1\). In \(\mathbb Z/N\mathbb Z\) define \[
 s=F^{-1},\qquad s^2=2,\qquad h=3+2s,\qquad q_0=-F.
\] The identity \((1+s)^K=Q_K+P_Ks=1\) gives \(h^K=1\). Moreover, \(h-1=2(1+s)\) is a unit since \((1+s)(1-s)=-1\). Modulo each prime \(\ell\mid N\), the order of \(h\) is exactly \(K\): it divides the prime \(K\) and is not one. Hence \(\ell\equiv1\pmod K\). In particular, \[
 1-h^j\text{ is a unit for }1\leq j<K,
 \qquad 1+h^j\text{ is a unit for every }j. \tag{PB.1}
\] The latter assertion uses odd \(K\). Thus \(H=\langle h\rangle\) has order \(K\) and acts freely on every nonzero residue, not just on the units. The signed orbits have size \(2K\), and \[
 n=\frac{N-1}{2K}
\] is an integer. The moduli need not be prime: for example \(Q_{11}=8119=23\cdot353\).

\subsection{Operators and identification with a fixed observable}\label{pell_blocks:operators-and-identification-with-a-fixed-observable}

On \(\mathbb C^{\mathbb Z/N\mathbb Z}\), with counting-measure inner product, write \[
 W(a,b)f(x)=e_N(a(x-b/2))f(x-b),\qquad
 S_u f(x)=f(u^{-1}x),\qquad Rf(x)=f(-x).
\] Here \(e_N(t)=e^{2\pi it/N}\), and divisions in exponents mean inverses modulo N. The symmetrized observables and spectral projectors preserve the odd subspace \(R=-1\). Individual Weyl translations do not, in general, preserve parity. For a character \(\eta\) of H, let \(E_{\eta}\) be the \(\eta\) eigenspace of \(S_{h}\) in that odd subspace, and let \(P_{\eta}\) denote its orthogonal projection. Its dimension is n.

Set \[
 B_q=W(q,1)+W(q,1)^*,\qquad
 A_{q,\eta}=P_\eta B_qP_\eta|_{E_\eta},\qquad
 X_{q,\eta}=\sqrt{K/2}\,A_{q,\eta}. \tag{PB.2}
\] In the theorem \(q=q_0\) is deterministic and is a unit modulo N.

This is a compression of a fixed Fourier observable of the fixed cat map, not merely a modulus-dependent choice of an unrelated classical map. Indeed, over \(\mathbb Z/N\mathbb Z\) put \[
 M=\begin{pmatrix}-1/(2s)&1\\-1/2&-s\end{pmatrix},\qquad
 T=M\operatorname{diag}(h,h^{-1})M^{-1}. \tag{PB.3}
\] The matrix \(M\) has determinant one. Define \(C_tf(x)=e_N(tx^2/2)f(x)\) and the unitary Fourier transform \(F_N\) with kernel \(e_N(xy)/\sqrt N\). Direct substitution gives their Weyl covariance matrices \[
 C_t:\begin{pmatrix}1&t\\0&1\end{pmatrix},\quad
 S_u:\begin{pmatrix}u^{-1}&0\\0&u\end{pmatrix},\quad
 F_N:\begin{pmatrix}0&1\\-1&0\end{pmatrix}.
\] Therefore \[
 V=C_{s^{-1}}S_{1/2}F_NC_{2s},\qquad U_N(T)=VS_h^{-1}V^* \tag{PB.4}
\] satisfies the exact finite Egorov identity \[
 U_N(T)W(v)U_N(T)^*=W(Tv).
\] These formulas also prove unitarity without invoking a prime-modulus-only Weil construction. They preserve parity. A global phase convention for the cat propagator does not affect any compression or relative eigenphase below.

Take the fixed observable \[
 B=W(1,0)+W(-1,0),\qquad Bf(x)=2\cos(2\pi x/N)f(x). \tag{PB.5}
\] Under conjugation by \(V^*\), its Weyl index is \(M^{-1}(1,0)=(-s,1/2)\). Conjugating further by \(S_2\), which commutes with \(S_h\), gives \((-s/2,1)=(-F,1)\). Thus the compressions of (PB.5) to every odd cat eigenspace are unitarily equivalent to \(A_{q_0,\eta}\), with a harmless permutation of sector labels.

\subsection{Main theorem}\label{pell_blocks:main-theorem}

\textbf{Theorem 1 (deterministic global Gaussian blocks).} Let K run through odd primes, with N,\(Q_{K}\),F,h,\(q_{0}\) as above. For every fixed integer \(m\geq1\), \[
 \sup_\eta\left|
 \frac1n\operatorname{Tr}X_{q_0,\eta}^{m}-\gamma_m
 \right|\longrightarrow0,
 \qquad
 \gamma_{2r}=(2r-1)!!,\quad\gamma_{2r+1}=0. \tag{PB.6}
\] Consequently the empirical spectral measures converge weakly to the standard normal distribution, uniformly over all K odd-sector cat eigenspaces. There is no random q in this theorem. For fixed even m the proof gives an \(O_{m}(\log K/K)\) error, apart from an exponentially small polynomial-in-K term. For fixed odd m the main word contribution is eventually absent.

The name Gaussian here describes the limiting \textbf{eigenvalue density}. It does not assert that the matrices form a Gaussian orthogonal ensemble.

\subsection{Exact traces over the composite ring}\label{pell_blocks:exact-traces-over-the-composite-ring}

The finite Weyl product remains \[
 W(a,b)W(c,d)=e_N((ad-bc)/2)W(a+c,b+d).
\] In an expansion of \((B_qP_\eta)^m\), let \(\epsilon_j\in\{1,-1\}\) and \(h_j\in H\). Starting with \((g,a,b,c)=(1,0,0,0)\), use \[
 a'=a+\epsilon/g,\quad b'=b+\epsilon g,\quad
 c'=c+\epsilon(ag-b/g)/2,\quad g'=gh. \tag{PB.7}
\] The word is \(e_N(qc)W(qa,b)S_g\), with the additional sector weight \(\overline{\eta(g)}\). For \(g\ne1\), (PB.1) gives a unique affine fixed point, so the trace has modulus one and phase \(qab(1+g)/(2(1-g))\). For \(g=1\), since q is a unit, the trace is N if \(a=b=0\) and zero otherwise. The parity insertion replaces g by \(-g\) and always has a unique fixed point by (PB.1).

For \(g=1\) write \(t_1=1\) and \(t_j=h_1\cdots h_{j-1}\). The remaining labels are arbitrary in H, and the final \(h_m\) is determined. The main term is independent of \(\eta\): \[
 \mu_{m,\eta}(q)=\frac{N}{N-1}
 \frac1{2^{m/2}K^{m/2-1}}
 \sum_{\substack{t_1=1,\ t_2,\ldots,t_m\in H;\ \epsilon\\
                  \sum\epsilon_jt_j=\sum\epsilon_j/t_j=0}}
 e_N(q\Omega)+\mathcal R_{m,\eta}(q), \tag{PB.8}
\] \[
 \Omega=\frac12\sum_{i<j}\epsilon_i\epsilon_j
                 (t_j/t_i-t_i/t_j),\qquad
 |\mathcal R_{m,\eta}(q)|\leq
 \frac{2^{m/2+1}K^{m/2+1}}{N-1}. \tag{PB.9}
\] This is an exact identity plus an elementary absolute bound. No field-only root-counting theorem or random-entry hypothesis is used.

\subsection{Lifting reciprocal relations into a quadratic field}\label{pell_blocks:lifting-reciprocal-relations-into-a-quadratic-field}

A tuple in (PB.8) gives \(P(Z)=\sum_j\epsilon_j Z^{e_j}\), where \(t_j=h^{e_j}\). Cancel opposite signs at equal labels first. If \(P\) is nonzero, its support has at most \(m\) points on \(\mathbb Z/K\mathbb Z\). Cut at a largest gap and rotate all labels. The resulting polynomial has coefficient \(\ell^1\)-norm at most \(m\) and degree \[
 L\leq K-\lceil K/m\rceil. \tag{PB.10}
\] Both reciprocal congruences survive this rotation.

Divide \(P\) by the monic polynomial \(Z^2-6Z+1\), obtaining an integer remainder \(aZ+b\). Since \(P(h)=P(h^{-1})=0\) and \(h-h^{-1}=4s\) is a unit modulo \(N\), both \(a\) and \(b\) are divisible by \(N\). Hence, for \(\lambda=3+2\sqrt2\), \[
 N^2\mid \operatorname{Norm}_{\mathbb Q(\sqrt2)/\mathbb Q}P(\lambda).
\] If this norm is nonzero, its absolute value is at most \[
 |P(\lambda)P(\lambda^{-1})|\leq m^2\lambda^L. \tag{PB.11}
\] But \(N^2\sim\lambda^K/4\), so (PB.10)--(PB.11) are incompatible for all sufficiently large K at each fixed m. Thus the rotated P satisfies \[
 P(\lambda)=0. \tag{PB.12}
\] An elementary sufficient numerical cutoff is \(\lambda^{K/m}>16m^2\). A sharper entirely rational check uses L in (PB.10) and \[
 m^2(2Q_{2L}+3P_{2L})<2N^2, \tag{PB.13}
\] since \(\sqrt2<3/2\). The implementation uses (PB.13), not an uncertified decimal comparison.

Crucially, (PB.12) does \textbf{not} imply coefficientwise balance. The structural relation \(h+h^{-1}-6=0\) is a nonbalanced reciprocal relation of length eight for every modulus in this family. An exclusion-of-all-relations programme would therefore be false even here.

\subsection{A bounded-cluster lemma and counting, instead of exclusion}\label{pell_blocks:a-bounded-cluster-lemma-and-counting-instead-of-exclusion}

\textbf{Lemma 2.} Let \(P\) be an integer polynomial of coefficient \(\ell^1\)-norm at most \(m\), with \(P(\lambda)=0\). Split its support at a gap \(d>2\log_\lambda m\). Then both pieces vanish separately at \(\lambda\).

\textbf{Proof.} Write \(P=P_{\rm low}+Z^{a+d}P_{\rm high}\), with \(\deg P_{\rm low}\leq a\). If \(P_{\rm high}(\lambda)\ne0\), its algebraic-integer norm is a nonzero integer. Since \(|P_{\rm high}(\lambda^{-1})|\leq m\), we obtain \(|P_{\rm high}(\lambda)|\geq1/m\). Cancellation with the low part would imply \(\lambda^{a+d}/m\leq m\lambda^a\), a contradiction to the gap assumption. Hence \(P_{\rm high}(\lambda)=0\), and then \(P_{\rm low}(\lambda)=0\). An empty low part is included. \(\square\)

Choose \[
 B_m=\lceil2\log_\lambda m\rceil+1,\qquad D_m=(m-1)B_m.
\] Repeatedly applying the lemma partitions a nonzero relation into components of span at most \(D_m\). Every nonzero component contains at least three signed occurrences; a two-term relation between positive distinct powers of \(\lambda\) is impossible. Reinsert the cancelled opposite-sign pairs as two-occurrence components of span zero.

If an m-tuple is not coefficientwise balanced, it therefore has at most \(\lfloor(m-1)/2\rfloor\) components. For each partition of the m positions, each sign pattern, and each bounded offset pattern, a component has one translation parameter in \(\mathbb Z/K\mathbb Z\). The component containing \(t_1=1\) has that parameter fixed. It follows that the number of unbalanced tuples in (PB.8) is at most \[
 C_m K^{\lfloor(m-1)/2\rfloor-1}. \tag{PB.14}
\] The constant depends only on \(m\). For example, \(C_m=(2m(2D_m+1))^m\) suffices by overcounting partitions, signs and offsets. Cyclic rotation introduces no further factor \(K\): the count is over cyclic component translations, with the component containing position 1 pinned.

For even \(m=2r\), (PB.14), divided by the normalization in (PB.8), is \(O_m(1/K)\), uniformly in q and \(\eta\). Odd m has an additional simplification. A polynomial vanishing at \(\lambda\) is divisible in \(\mathbb Z[Z]\) by \(Z^2-6Z+1\), so \(P(1)\) is divisible by \(-4\). Yet the parity of \(P(1)\) equals the parity of the number of signed occurrences. Consequently odd-m reciprocal tuples do not exist once the lifting argument applies.

\subsection{All pairings survive for the fixed observable}\label{pell_blocks:all-pairings-survive-for-the-fixed-observable}

For a balanced \(2r\)-tuple, match equal labels with opposite signs. Distinct pair labels give a unique matching; repeated labels contribute \(O_r(1/K)\) normalized overcount. For a matching \(\pi\), orient each pair by \(\sigma_i\in\{1,-1\}\), and let its label be \(x_i=h^{e_i}\), with \(x_1=1\). The exact phase is \[
 F_{\pi,\sigma}(x)=
 \sum_{\substack{i<j\\\text{pairs }i,j\text{ cross}}}
 \sigma_i\sigma_j(x_j/x_i-x_i/x_j). \tag{PB.15}
\] Unlike a q-average, the deterministic value \(q_0\) does not suppress crossings. Let \(\delta\) be the representative of d modulo K in \([-(K-1)/2,(K-1)/2]\). The Pell recurrence gives \[
 q_0(h^d-h^{-d})\equiv-2P_{2\delta}\pmod N. \tag{PB.16}
\] For \(|\delta|\leq(K-1)/2-L\), \[
 |P_{2\delta}|/N\leq C\lambda^{-L}. \tag{PB.17}
\] The implicit constant is absolute. A difference of two independent uniform pair labels is uniform on \(\mathbb Z/K\mathbb Z\), also when one label is fixed. For each fixed pairing a union bound shows that only \(O_r(L/K)\) assignments have a crossing difference in the excluded central window. On all other assignments, (PB.15)--(PB.17) give phase \(1+O_r(\lambda^{-L})\). Taking \(L=\lceil\log_\lambda K\rceil\) proves \[
 K^{-(r-1)}\sum_{x_2,\ldots,x_r\in H}
 e_N(q_0 F_{\pi,\sigma}(x))
 =1+O_r(\log K/K). \tag{PB.18}
\] There are \((2r-1)!!\) matchings, not just the Catalan number of noncrossing matchings. Equations (PB.8)--(PB.9), (PB.14), and (PB.18) prove (PB.6) and its stated even moment error. Odd moments follow from the parity argument and (PB.9).

Bounded second moments give tightness, and higher even moments give uniform integrability. The Gaussian moment sequence is determinate: its even moments imply a finite exponential-square moment for every exponent below \(1/2\), by monotone expansion, and hence determine the characteristic function \(e^{-t^2/2}\). The usual subsequence argument therefore gives weak convergence. Because all bounds are uniform in \(\eta\), selecting a worst sector on any hypothetically nonconvergent subsequence proves the claimed uniformity.

\subsection{An explicit fourth-moment certificate}\label{pell_blocks:an-explicit-fourth-moment-certificate}

For \(m=4\), a nonzero polynomial of coefficient \(\ell^1\)-norm at most \(4\) cannot vanish at \(\lambda>3\): its leading term dominates all lower terms. Thus, once (PB.13) holds, all four-letter reciprocal tuples are balanced. Exact pairing overcount is six, giving uniformly in \(\eta\), \[
 \left|\mu_{4,\eta}-\frac{N}{N-1}
 \left(2+S_K-\frac{3}{2K}\right)\right|
 \leq\frac{8K^3}{N-1}, \tag{PB.19}
\] \[
 S_K=\frac1K\left(1+2\sum_{d=1}^{(K-1)/2}
                    \cos(4\pi P_{2d}/N)\right). \tag{PB.20}
\] These are finite-size bounds, not fitted matrix histograms. The maintained companion reports high-precision cosines and an analytic series-tail bound, separately from floating-point rounding. It includes the full remainder in (PB.19), but does not provide rational interval enclosures.

There is also a nontrivial explicit asymptotic correction: \[
 K(3-\mu_{4,\eta})\longrightarrow
 \frac32+2\sum_{j=0}^\infty
 \left[1-\cos\left((4-2\sqrt2)\pi\lambda^{-j}\right)\right]. \tag{PB.21}
\] To see this, reindex (PB.20) by \(d=(K-1)/2-j\). Then \(P_{K-1-2j}/N\to\alpha^{-1-2j}/\sqrt2\). The inequality \(1-\cos x\leq x^2/2\) gives a uniform summable geometric majorant. The remainder in (PB.19), even after multiplication by \(K\), is exponentially small. This proves (PB.21), uniformly in the sector.

\subsection{A finite-perturbation consequence for actual eigenphases}\label{pell_blocks:a-finite-perturbation-consequence-for-actual-eigenphases}

The following elementary comparison explains precisely how far the global observable theorem transfers to a propagator.

\textbf{Lemma 3.} Let \(U_0\) be unitary, let \(P\) be its spectral projection at \(1\), and let \(\delta>0\) be the distance from \(1\) to the other eigenvalues. Let \(B=B^*\), \(\|B\|\leq b\), and \(U_e=e^{ieB}U_0\). If \(0<b|e|\leq\delta/16\), the cluster near \(1\) has rank \(\operatorname{rank}P\). Write its principal eigenphases as \(\theta_j\), and the eigenvalues of \(A=PBP\) on \(\operatorname{ran}P\) as \(a_j\). Their empirical measures satisfy \[
 W_2\left(\frac1n\sum_j\delta_{\theta_j/e},
          \frac1n\sum_j\delta_{a_j}\right)
 \leq\frac{13\pi b^2|e|}{\delta}. \tag{PB.22}
\] Here \(W_2\) is quadratic Wasserstein distance between equal-weight real spectra.

\textbf{Proof.} Put \(\Delta=U_e-U_0\) and \(d=\|\Delta\|\leq b|e|\). Resolvent integration on \(|z-1|=\delta/2\) gives \(\|P_e-P\|=\eta\leq4d/\delta\leq1/4\). The new projection is orthogonal. Set \(H=(PP_eP)^{1/2}\) on \(\operatorname{ran}P\) and \(J=P_ePH^{-1}\). Then \(J\) is an isometry onto the perturbed cluster, \(PJ=H\), \(\|(I-P)J\|\leq\eta\), and \(\|J-P\|\leq2\eta\). Let \(V=J^*U_eJ\). Spectral inclusion for the normal matrix \(U_0\) gives \(\|V-I\|\leq d\): every cluster eigenvalue is within \(d\) of \(1\), since all other old eigenvalues are farther than \(\delta/2\).

Writing \(Z=(I-P)J\) and using invariance of the new cluster gives \((U_0-I)Z=Z(V-I)-(I-P)\Delta J\). Hence \(\|Z^*(U_0-I)Z\|\leq\eta(1+\eta)d\). Also \(\|J^*\Delta J-P\Delta P\|\leq4\eta d\). Therefore \(\|V-I-P\Delta P\|\leq24d^2/\delta\). Taylor's integral remainder for the two unitary exponentials, and \(\delta\leq2\), yield \[
 \|V-e^{i e A}\|\leq26 e^2b^2/\delta. \tag{PB.23}
\] For two unitary matrices, the normalized squared Frobenius difference is a doubly stochastic average of squared eigenvalue distances, using the squared absolute overlaps of their eigenbases. Decomposing that matrix into permutation matrices gives a matching whose mean squared chord distance is at most the squared operator norm in (PB.23). Such a decomposition follows by repeatedly selecting a positive perfect matching (Hall's condition follows from row and column sums), subtracting its smallest entry, and iterating. All phases lie in one short arc. For angular differences at most \(\pi\), arc distance is at most \(\pi/2\) times chord distance. Dividing by \(|e|\) proves (PB.22). \(\square\)

For the cat propagator, rotate any of its K clusters to 1. Its gap is \[
 \delta_K=2\sin(\pi/K)\geq4/K,\qquad b=2.
\] Thus the empirical law of the rescaled relative eigenphases \(\sqrt{K/2}\,\theta_j/e\) differs in \(W_2\) from the corresponding observable block by at most \[
 \frac{13\pi}{\sqrt2}|e|K^{3/2}. \tag{PB.24}
\] \textbf{Corollary 4.} For any nonzero perturbation sequence \(e=e_K\) satisfying \(|e_K|K^{3/2}\to0\), every odd-sector eigenphase cluster of \(e^{ie_KB}U_N(T)\), rescaled as above, converges globally to the standard normal distribution, uniformly over the K clusters. The symmetrically kicked propagator \(e^{ie_KB/2}U_N(T)e^{ie_KB/2}\) is isospectral and has the same conclusion.

This does not exchange the small-perturbation and semiclassical limits without control: (PB.24) is the control, and its hypothesis is part of the claim. Fixed nonzero perturbations, microscopic within-cluster correlations, and full BGS universality are not established.
